\clearpage\thispagestyle{empty}
\pdfbookmark[0]{Quellen zur Besprechung der Connoissance des tems}{connoissance1808-sources}
\begin{center}\textbf{Quellen zur Besprechung der Connoissance des tems}\end{center}
\noindent Diese KI-Übertragung ins heutige Deutsch folgt der Besprechung von Gauß auf S.~492--496 der \href{https://gdz.sub.uni-goettingen.de/id/PPN236007467}{\textit{Werke}, Bd.~VI} (Göttingen, 1874). Sie erschien ohne Namensnennung unter der Überschrift \emph{Paris} in den \textit{Göttingischen gelehrten Anzeigen}, Stück~56 und~57, 7. April 1808, S.~\href{https://api.digitale-sammlungen.de/iiif/image/v2/bsb10538673_00559/full/1800,/0/default.jpg}{553}--\href{https://api.digitale-sammlungen.de/iiif/image/v2/bsb10538673_00568/full/1800,/0/default.jpg}{562}. Das verlinkte Exemplar stammt aus der Bayerischen Staatsbibliothek.
\par\smallskip\noindent Der französische Titel bedeutet \emph{Kenntnis der Zeiten oder der Himmelsbewegungen zum Gebrauch für Astronomen und Seefahrer}. Herausgeber war das Bureau des longitudes. Die Werkausgabe datiert die bei Lalande erwähnte Sonnenfinsternis auf den 1. Februar 1804. Im Erstdruck steht auf \href{https://api.digitale-sammlungen.de/iiif/image/v2/bsb10538673_00561/full/1800,/0/default.jpg}{S.~555} der 11. Februar 1804. Die Übertragung folgt der Werkausgabe. KI-Übertragung und KI-Satz.
\clearpage
